\chapter{Nix Workflow}

\section{flake.nix}

The repository includes a \texttt{flake.nix} that declares all TeX Live
packages and development dependencies:

\begin{minted}{nix}
{
  inputs = {
    nixpkgs.url = "github:NixOS/nixpkgs/nixos-unstable";
    flake-utils.url = "github:numtide/flake-utils";
  };

  outputs = { nixpkgs, flake-utils, ... }:
    flake-utils.lib.eachDefaultSystem (system:
      let pkgs = nixpkgs.legacyPackages.${system};
      in {
        devShells.default = pkgs.mkShell {
          packages = with pkgs; [
            (texlive.combine {
              inherit (texlive) scheme-full ...;
            })
            python3
            lean4
            lake
          ];
        };
      }
    );
}
\end{minted}

\section{Dev Shell}

Activate the development environment with:

\begin{minted}{bash}
nix develop .#
\end{minted}

This provides a shell with:
\begin{itemize}
  \item TeX Live (full scheme with all OmniLaTeX dependencies).
  \item Python~3 (for \texttt{build.py} and \texttt{pygments}).
  \item Lean~4 and Lake (for proof verification).
\end{itemize}

\section{Lean 4}

Proof-heavy documents can use Lean~4 for formal verification:

\begin{minted}{bash}
lake build
\end{minted}

The dev shell includes \texttt{lean4} and \texttt{lake}, so no separate
installation is needed.

\section{Reproducibility}

\texttt{flake.lock} pins every dependency to an exact revision, ensuring
bit-for-bit reproducible builds across machines.  Commit
\texttt{flake.lock} alongside your source to guarantee that collaborators
and CI use identical package versions.

\section{Multi-Platform}

The flake supports multiple systems:

\begin{itemize}
  \item \texttt{x86\_64-linux} --- primary target, fully tested.
  \item \texttt{aarch64-darwin} --- Apple Silicon (macOS).
  \item \texttt{x86\_64-darwin} --- Intel macOS.
\end{itemize}

\section{CI Integration}

GitHub Actions workflows use Nix for reproducible CI:

\begin{minted}{yaml}
- uses: DeterminateSystems/nix-installer-action@main
- uses: DeterminateSystems/magic-nix-cache-action@main
- run: nix develop .# --command bash -c \
      "BUILD_MODE=prod python build.py"
\end{minted}

The \texttt{magic-nix-cache-action} enables binary cache sharing between CI
runs, reducing build times significantly.
